\section{Introduction}

% ¶1: Context — agents are powerful and widely deployed through harnesses
AI agents increasingly operate as long-running actors with access to
tools, files, shells, browsers, package managers, and external
APIs~\cite{debenedetti2024agentdojo,zhan2024injecagent,
chennabasappa2025llamafirewall,agentsight}. Coding agents are a
prominent example~\cite{yang2024sweagent,wang2025openhands}. As agents gain autonomy,
they require \emph{AI agent harnesses}: software around the model that
improves agent performance while enforcing behavioral constraints for
safety, security, and compliance~\cite{langchain-harness,fowler-harness}.

% ¶2: Policy engines; instruction files are behavioral policy
A central component of such a harness is a policy engine: the
part that enforces behavioral constraints over the agent's concrete
actions. In practice, projects encode many such constraints as
behavioral policies in instruction files
(\texttt{CLAUDE.md},
\texttt{AGENTS.md})~\cite{chatlatanagulchai2025manifests,lulla2026agentsmd}, so
harnesses can provide these instructions text to the model directly.
Because LLMs itself comply with instructions
probabilistically~\cite{liu2024lost,jiang2024followbench,qi2025agentif}
and may violate them through planning errors, 
prompt-injection drift, and opaque tool or scripts
behavior~\cite{greshake2023indirectprompt,zhan2024injecagent,debenedetti2024agentdojo},
harnesses also require deterministic policy engines to enforce compliance at 
runtime~\cite{rebedea2023nemo,wang2025agentspec,shi2025progent,
costa2025fides,debenedetti2025camel}.

% ¶4: Semantic gap — existing deterministic mechanisms also fail
Many of these policies constrain system-level actions, yet existing policy engines in agent harnesses leave a \emph{semantic gap}:
\emph{AI agent intent and directives} are expressed in natural language, but enforcement must decide over
\emph{system actions}.
Our
empirical study of 64~projects (\S\ref{sec:empirical}) finds that 64\%
of instruction-file statements are behavioral directives, and 83\% of
those involve system-level behavior: commands executed, files accessed,
and network connections made.
Tool-call based policy
systems in the Agent runtime like claude code, codex with LLM guardrails or regex-like checkers~\cite{rebedea2023nemo,xiang2025guardagent,wang2025agentspec,
shi2025progent,costa2025fides,debenedetti2025camel}
lose track of system actions when agents shell out or spawn subprocesses: the
same \texttt{git~push} can be reached through a tool call, a
\texttt{bash~-c} wrapper, or a compiled binary created by previous session
; and at the repository level, 81\% of projects
contain cross-event directives that a single tool-level interception cannot
enforce, e.g. run test before commit.

Static policy is insufficient, 
the policy engine needs to define and evaluate policies at runtime, 
and allowing the agent to adapt them based on project or task context.
We found that 74\% of system-level directives require project or
task context that is absent from low-level OS events. OS-level
enforcement system and sandbox for traditional software
~\cite{seccomp-bpf,apparmor,landlock,ciliumtetragon,flume,
pasquier2018camquery} can observe all system actions, but expect pre-defined
static policy and return only system events without semantic context, making the 
agent unable to follow the instructions or adapt the rules.

% ¶5: Insight — programmable policy engine below the tool layer
We argue that an agent-harness policy engine should be
programmable define by the agent itself and take effect at OS level: 
it should let the agent use live project and
task context to enforce its own actions and adapt at runtime. 
This creates three challenges. First, policy specification must be \emph{expressive} enough to
capture directive intent yet \emph{concrete} enough to compile to
deterministic kernel checks. Second, enforcement must track
policy-relevant state across tools, subprocesses, files, and network
endpoints without imposing prohibitive per-syscall overhead. Third,
programmability must not become an authority bypass: agents must be able
to add context-specific constraints, but neither they nor their
sub-agents may weaken inherited policy.

% ¶6: System — \sys{}
\sys{} is a programmable OS-level runtime policy enforcement system for
AI agent harnesses. Agents express behavioral policies in a compact
domain-specific language (DSL) generated from natural-language project
directives; \sys{} compiles them into labeled information-flow rules and
loads them into an eBPF/BPF-LSM
(Linux Security Modules) enforcement engine. Named labels propagate across processes, files, and network
endpoints at kernel hooks, encoding execution history as checkable state.
When a labeled event matches a rule, \sys{} returns semantic
feedback: which rule matched, why, and what compliant alternative
satisfies it. Each rule independently selects one of three effects:
notify (observe), block (pre-operation denial), or kill (post-operation
termination). A layered authority model ensures that higher-authority
rules (e.g., platform-level ``do not leak secrets'') cannot be weakened by
agent-authored policy.

We make three contributions:

\begin{enumerate}
\item An \textbf{empirical study} of instruction files from 64~projects,
showing that cross-event system-level directives are common (81\% of
repositories) and that 74\% of system-level directives require project or
task context absent from low-level OS events (\S\ref{sec:empirical}).

\item \textbf{\sys{}}, a programmable OS-level runtime policy enforcement
system for AI agent harnesses. \sys{} compiles behavioral policies into eBPF/BPF-LSM labeled
information-flow rules, enforces them across process, file, and network
boundaries, and returns semantic feedback that explains rule matches in
terms of AI agent intent and concrete remediation
(\S\ref{sec:design}--\ref{sec:implementation}).

\item An \textbf{evaluation} across directive-derived rules,
system workloads, and a 20-task OctoBench subset. On a 190-trace
decision-compliance benchmark sampled from 38~directives, \sys{} reaches
a 75.8\% decision-compliance rate (DCR), 27.4 percentage points higher than prompt-filter and
30.5 points higher than tool-regex, while adding less than 2\% end-to-end overhead on
agent workloads
(\S\ref{sec:evaluation}).
\end{enumerate}
